\documentclass[aspectratio=169,10pt]{beamer}
\usepackage{microtype}
\usepackage{amsmath,amssymb,mathtools}

\usetheme{metropolis}
\usefonttheme{professionalfonts}
\setbeamertemplate{navigation symbols}{}
\setbeamertemplate{theorems}[numbered]

% ── Palette ────────────────────────────────────────────────────────────────
\definecolor{cpnavy}{HTML}{1A3047}   % deep navy — frametitle, headings
\definecolor{cpgold}{HTML}{C9943A}   % warm gold — accents, alerts
\definecolor{cpstone}{HTML}{F5F3EF}  % warm off-white — body bg
\definecolor{cpfog}{HTML}{EAE6DF}    % light warm gray — block backgrounds
\definecolor{cpmid}{HTML}{5C6B7A}    % medium slate — secondary text

\setbeamercolor{normal text}{fg=cpnavy,bg=cpstone}
\setbeamercolor{frametitle}{bg=cpnavy,fg=white}
\setbeamercolor{title separator}{fg=cpgold}
\setbeamercolor{progress bar}{fg=cpgold,bg=cpnavy!30}
\setbeamercolor{alerted text}{fg=cpgold}
\setbeamercolor{block title}{bg=cpnavy,fg=white}
\setbeamercolor{block body}{bg=cpfog,fg=cpnavy}
\setbeamercolor{block title alerted}{bg=cpgold,fg=white}
\setbeamercolor{block body alerted}{bg=cpfog,fg=cpnavy}
\setbeamercolor{itemize item}{fg=cpgold}
\setbeamercolor{enumerate item}{fg=cpgold}
\setbeamercolor{itemize subitem}{fg=cpmid}

\setbeamerfont{block title}{size=\small,series=\bfseries}
\setbeamerfont{frametitle}{size=\large,series=\bfseries}

\newcommand{\RR}{\mathbb{R}}
\newcommand{\CP}{\mathcal{C}}
\newcommand{\eps}{\varepsilon}
\newcommand{\norm}[1]{\left\lVert #1\right\rVert}

\title{\textbf{CP Generator}}
\subtitle{Generation, Diagnostics, and Animation of Origami Crease Patterns}
\author{Allan Pan}
\date{\today}

\begin{document}

\maketitle

% ────────────────────────────────────────────────────────────────────────────
\begin{frame}{What this project is}
\begin{columns}[T,onlytextwidth]
\column{0.52\textwidth}
A desktop app that generates, diagnoses, and animates origami crease patterns
on a square sheet — entirely from scratch.
\medskip

\textcolor{cpmid}{\small \texttt{src/cp\_generator/core.py} · \texttt{src/cp\_generator/app.py} · \texttt{src/cp\_generator/fold\_sim.py}}
\medskip

\begin{block}{Five linked algorithms}
\begin{enumerate}\small
  \item Build a planar candidate graph.
  \item Repair parity \& optimize for Kawasaki.
  \item Assign mountains/valleys.
  \item Extract faces \& animate the fold.
  \item Diagnose what is only local, what is globally consistent, and what is still heuristic.
\end{enumerate}
\end{block}

\column{0.44\textwidth}
\textbf{Key objects}
\begin{itemize}\small
  \item Sheet \([0,s]^2\); graph \(G=(V,E)\), positions \(X_v\in\RR^2\).
  \item Partial crease labeling \(\tau:E\to\{U,M,V\}\).
  \item Complete MV assignment \(\sigma:E\to\{-1,+1\}\), with \(+1=M\) and \(-1=V\).
  \item \alert{Sector angle}: angle between consecutive creases at a vertex.
  \item \alert{Half-edge}: oriented edge — determines which face lies to its left.
\end{itemize}
\end{columns}
\end{frame}

% ────────────────────────────────────────────────────────────────────────────
\begin{frame}{Algorithm I — Planar candidate graph}
\textbf{Goal.} A planar graph dense enough for interesting folds, regular
enough for stable local angle computation.
\medskip

\begin{columns}[T,onlytextwidth]
\column{0.54\textwidth}
\begin{block}{Construction}
\begin{enumerate}\small
  \item Sample interior points; add four corners.
  \item Snap near-boundary points to the boundary.
  \item Compute the Delaunay triangulation.
  \item Use its edges as the initial crease graph.
\end{enumerate}
\end{block}

\column{0.42\textwidth}
\textbf{Why Delaunay?}\\[3pt]
\small No other input point lies inside the circumcircle of any triangle —
avoids skinny triangles and gives clean local angle data.
\end{columns}
\end{frame}

% ────────────────────────────────────────────────────────────────────────────
\begin{frame}{Algorithm II — Parity repair \& Kawasaki optimization}
\begin{columns}[T,onlytextwidth]
\column{0.48\textwidth}
\begin{alertblock}{Necessary condition}
Every interior vertex must have \alert{even degree}:
\(d(v)\equiv 0\pmod 2.\)
\end{alertblock}
\smallskip
\small\textbf{Repair.} Pair odd-degree vertices; delete the shortest path
between each pair. \textcolor{cpmid}{(Heuristic — not globally optimal.)}

\column{0.48\textwidth}
\begin{block}{Kawasaki's theorem (local)}
Sector angles \(\alpha_1,\dots,\alpha_{2m}\) must satisfy
\[
\alpha_1+\alpha_3+\cdots = \alpha_2+\alpha_4+\cdots = \pi.
\]
\end{block}
\end{columns}

\medskip
\textbf{Geometry repair.} Let \(D(X)=\sum_{v} w_v\,\norm{X_v-X_{0,v}}^2\).
Move vertices as little as possible, subject to the Kawasaki constraints:
\[
\min_{X}\;D(X)+\frac{\alpha}{2}\,D(X)^{2},
\]
Boundary weights \(w_v\) are large; interior weights are small. Solved with SLSQP.
\end{frame}

% ────────────────────────────────────────────────────────────────────────────
\begin{frame}{Algorithm III — Local constraints: Maekawa \& Bern--Hayes}
\begin{columns}[T,onlytextwidth]
\column{0.46\textwidth}
\begin{block}{Maekawa's theorem (local)}
At every interior vertex,
\[
M(v)-V(v)=\pm 2.
\]
Used as an \textbf{admissibility filter}: reject any assignment violating
this anywhere.
\end{block}

\column{0.50\textwidth}
\textbf{Bern--Hayes-style reduction}\\[4pt]
For each vertex, repeatedly \alert{remove the smallest sector}: its two
bounding creases are forced \emph{equal} or \emph{opposite} in sign.
\smallskip
\begin{enumerate}\small
  \item Sort creases around the vertex by angle.
  \item Find the smallest wedge; record its forced sign relation.
  \item Recurse until the vertex is fully resolved.
\end{enumerate}
\smallskip
\textcolor{cpmid}{\small Boundary vertices get temporary virtual rays to close the cyclic order.}
\end{columns}
\end{frame}

% ────────────────────────────────────────────────────────────────────────────
\begin{frame}{Algorithm III — Merging, exact search, \& interlacing}
Each vertex reduction yields sign chains:
\[
\sigma(e_i)=\sigma(e_j)\quad\text{or}\quad\sigma(e_i)=-\sigma(e_j).
\]
\begin{columns}[T,onlytextwidth]
\column{0.52\textwidth}
\textbf{Exact search}
\begin{itemize}\small
  \item Merge chains across shared creases into \(g\) free groups.
  \item Enumerate all \(2^g\) seed assignments; propagate implied signs.
  \item Keep only assignments that pass Maekawa at every vertex.
\end{itemize}
\textcolor{cpmid}{\small Undetermined creases keep label \(-1\); no sign is invented.}

\column{0.44\textwidth}
\begin{block}{Interlacing preference}
\small Among admissible assignments, prefer
\(M,V,M,V,\ldots\) over \(M,M,V,V,\ldots\); weight sharper wedges more.\\[4pt]
\alert{Does not affect correctness.}
\end{block}
\end{columns}
\end{frame}

% ────────────────────────────────────────────────────────────────────────────
\begin{frame}{Algorithm IV — Faces, reflection \& 3D animation}
\begin{columns}[T,onlytextwidth]
\column{0.48\textwidth}
\textbf{Face extraction}\\[4pt]
\small Trace half-edges in order; each bounded counterclockwise cycle defines one face.
Faces are then triangulated by ear clipping for rendering.

\medskip
\begin{block}{Reflection rule}
\small If two faces share crease \(\ell\), the child is the parent reflected
in \(\ell\). Propagated from a root face outward; if any cycle is inconsistent,
the method falls back to per-triangle propagation along a spanning tree.
\end{block}

\column{0.48\textwidth}
\textbf{Kinematic animation}\\[4pt]
Around crease axis \((p,q)\), with sign \(\sigma(e)\in\{\pm 1\}\) and fold
angle \(\theta\):
\[
T_{\text{child}} = R_{(p,q)}\!\bigl(\sigma(e)\,\theta\bigr)\,T_{\text{parent}}.
\]
\small Rigid plates connected by hinge creases — no elasticity, no collisions.
Near the fully folded state, small offsets along the surface normal keep
layer order visible.
\end{columns}
\end{frame}

% ────────────────────────────────────────────────────────────────────────────
\begin{frame}{Algorithm V — Diagnostic status map}
\footnotesize Let \(\CP=(G,X,\tau)\) be the current crease pattern, let \(V_{\mathrm{int}}\subseteq V\) denote the interior vertices, and let
\[
\kappa(v)=\left|\sum_{i\text{ odd}}\alpha_i(v)-\pi\right|
\]
be the Kawasaki residual at \(v\), where \(\alpha_1(v),\dots,\alpha_{2m(v)}(v)\) are the sector angles in cyclic order.

\begin{block}{Reported statuses}
\begin{enumerate}\footnotesize
  \item \textbf{Local status = pass} iff every \(v\in V_{\mathrm{int}}\) has even degree and \(\kappa(v)\le \eps\).
  \item \textbf{Assignment status = pass} iff \(\tau\) is complete and \(|M_\sigma(v)-V_\sigma(v)|=2\) for every \(v\in V_{\mathrm{int}}\).
  \item \textbf{Global status = pass} iff the current embedding has no nonincident crease crossings, every crease borders exactly two traced faces, and reflection propagation is cycle-consistent on the face-adjacency graph.
  \item \textbf{Exact-preview status = pass} iff the current geometry admits exact face reconstruction with no provisional signs and no mesh fallback.
\end{enumerate}
\end{block}

\begin{alertblock}{Non-certifying outcomes}
\footnotesize A warning is reported whenever exact reconstruction needs a reference-geometry reset, provisional completion of \(\tau\), or mesh fallback.
\end{alertblock}
\end{frame}

% ────────────────────────────────────────────────────────────────────────────
\begin{frame}{Algorithm VI — Implemented global-consistency test}
\footnotesize The current code does \emph{not} solve arbitrary global flat-foldability. It implements the following certifier on the current embedding \(X\):

\begin{block}{Implemented decision procedure}
\begin{enumerate}\footnotesize
  \item If nonincident fold segments cross, return \(\texttt{fail}\). If there are no folds, return \(\texttt{not\_run}\). If no MV assignment has been attempted, return \(\texttt{unknown}\).
  \item Augment the fold graph by the square boundary, and order neighbors at each vertex by polar angle.
  \item Trace bounded faces by the left-face rule \((u,v)\mapsto(v,w)\), where \(w\) is the predecessor of \(u\) in the cyclic order at \(v\). Reject repeated vertices or nonclosing walks.
  \item Require every fold edge to border exactly two traced faces. Build the face-adjacency graph, choose a maximal-area root face, and set \(T_{f_0}=I\).
  \item Across each crease \(\ell\), propagate \(T_g=R_{\ell}T_f\). On revisiting a face, compare the competing transforms: hard discrepancy gives \(\texttt{fail}\), soft discrepancy gives \(\texttt{warning}\).
  \item Resolve fold signs from \(\tau\). Any missing sign is filled provisionally by face-depth parity, which downgrades the outcome to \(\texttt{warning}\).
\end{enumerate}
\end{block}

\begin{alertblock}{Output interpretation}
\footnotesize A final \(\texttt{pass}\) means: no crossings, successful exact face tracing, two-sided incidence for every fold, cycle-consistent reflections within tolerance, and no provisional MV completion. This is the strongest global certificate currently implemented in the code.
\end{alertblock}
\end{frame}

% ────────────────────────────────────────────────────────────────────────────
\begin{frame}{Scope of the local certificates}
\begin{theorem}[Hull's local flat-foldability conditions]
If a single-vertex crease pattern flat-folds, then its degree is even, its sector angles satisfy Kawasaki's alternating-sum identity, and every complete MV assignment satisfies Maekawa's relation \(|M-V|=2\). \textcolor{cpmid}{\cite{hull_math_flat}}
\end{theorem}

\begin{proof}[Interpretation in this project]
\footnotesize The implemented \emph{local status} and \emph{assignment status} test these vertexwise necessary conditions, up to numerical tolerance \(\eps\). Bern--Hayes reduction is used only to derive local sign constraints and shrink the search over complete assignments \(\sigma\); it does not decide global flat-foldability for a multi-vertex pattern. Hence a pass certifies \emph{local admissibility}, not a global folding. \textcolor{cpmid}{\cite{hull_math_flat,bern_hayes}}
\end{proof}

\begin{block}{Why this distinction matters}
\footnotesize Hull's theorem is local, while Bern--Hayes shows that exact flat foldability becomes globally difficult once many vertices interact. The report therefore keeps ``locally admissible assignment'' and ``globally certified current geometry'' separate. \textcolor{cpmid}{\cite{bern_hayes}}
\end{block}
\end{frame}

% ────────────────────────────────────────────────────────────────────────────
\begin{frame}{Meaning of a global or exact-preview pass}
\begin{theorem}[Implemented exact-current-geometry certificate]
If the algorithm on the previous slide returns \emph{global status = pass} and \emph{exact-preview status = pass}, then the current embedding of \(\CP\) has a noncrossing augmented planar graph, a closed family of bounded traced faces, a two-face incidence relation on every fold edge, and a cycle-consistent family of flat reflections \(T_g=R_{\ell}T_f\), all without provisional completion of the MV data.
\end{theorem}

\begin{proof}[Why this theorem is faithful to the code]
\footnotesize
\begin{enumerate}\itemsep2pt
  \item `src/cp_generator/core.py` first rejects nonincident segment crossings.
  \item `fold\_sim.py` then traces bounded faces by cyclic half-edge traversal and aborts on repetition or nonclosure.
  \item The exact solver requires exactly two incident faces per fold and propagates transforms by explicit crease reflections.
  \item Cycle drift beyond tolerance aborts; smaller drift or provisional sign completion downgrades to warning rather than pass.
\end{enumerate}
Thus the theorem states exactly what the present code certifies. It is not yet a complete decision procedure for arbitrary multi-vertex flat-foldability, but it is consistent with the formal viewpoint that valid flat foldings require noncrossing geometry and coherent global compatibility data. \textcolor{cpmid}{\cite{hull_zakharevich,eppstein_flat}}
\end{proof}
\end{frame}

% ────────────────────────────────────────────────────────────────────────────
\begin{frame}{Certified today vs.\ unresolved}
\begin{columns}[T,onlytextwidth]
\column{0.46\textwidth}
\begin{block}{Currently certifiable}
\begin{itemize}\small
  \item Vertexwise even-degree and Kawasaki residual checks.
  \item Maekawa verification on a complete assignment \(\sigma\).
  \item Explicit detection of nonincident crease crossings.
  \item Exact face reconstruction on the current geometry.
\end{itemize}
\end{block}

\column{0.50\textwidth}
\begin{block}{Still non-certifying}
\begin{itemize}\small
  \item Parity repair is not globally optimal.
  \item Bern--Hayes reduction is local, not a full global classifier.
  \item Reference-geometry and mesh preview fallbacks remain heuristic.
  \item A true exact global checker is still needed for larger patterns.
\end{itemize}
\end{block}
\end{columns}

\begin{alertblock}{Best theoretical next step}
\footnotesize The most plausible upgrade path is a bounded-instance exact checker based on ply and cell-adjacency complexity, built on the face arrangement already computed by the exact preview solver. \textcolor{cpmid}{\cite{eppstein_flat}}
\end{alertblock}
\end{frame}

% ────────────────────────────────────────────────────────────────────────────
\begin{frame}[allowframebreaks]{References}
\footnotesize
\begin{thebibliography}{9}
\bibitem[BH96]{bern_hayes}
Marshall Bern and Barry Hayes.
\newblock \emph{The Complexity of Flat Origami}.
\newblock SODA, 1996.

\bibitem[Hul94]{hull_math_flat}
Thomas C. Hull.
\newblock \emph{On the Mathematics of Flat Origamis}.
\newblock Congressus Numerantium, 1994.

\bibitem[AE\!+\!15]{abel_rigid}
Zachary Abel, Jason Cantarella, Erik D. Demaine, David Eppstein, Thomas C. Hull, Jason S. Ku, Robert J. Lang, and Tomohiro Tachi.
\newblock \emph{Rigid Origami Vertices: Conditions and Forcing Sets}.
\newblock Journal of Computational Geometry, 2015.

\bibitem[Epp23]{eppstein_flat}
David Eppstein.
\newblock \emph{A Parameterized Algorithm for Flat Folding}.
\newblock arXiv:2306.11939, 2023.

\bibitem[HZ25]{hull_zakharevich}
Thomas C. Hull and Inna Zakharevich.
\newblock \emph{Flat Origami Is Turing Complete}.
\newblock arXiv:2309.07932, 2025 version consulted.

\bibitem[WS22]{walker_stankovic}
Andreas Walker and Tino Stankovi\'c.
\newblock \emph{Algorithmic Design of Origami Mechanisms and Tessellations}.
\newblock Communications Materials, 2022.
\end{thebibliography}
\end{frame}

\end{document}
